\begin{lemma}[Limit passage in the weak $p$-Laplace equation]
Let $1<p<2$, let $\Omega\subseteq\mathrm{Plane}$, let $u_n\colon\mathrm{Plane}\to\mathbb{R}$ and
$g_n\colon\mathrm{Plane}\to\mathrm{Plane}$ for $n\in\mathbb{N}$, and let
$U\colon\mathrm{Plane}\to\mathbb{R}$ and $G\colon\mathrm{Plane}\to\mathrm{Plane}$.  Assume:
\begin{enumerate}
\item for every $n\in\mathbb{N}$, the function $u_n$ is $p$-harmonic in $\Omega$ with weak
gradient $g_n$;
\item for every compact $K\subseteq\Omega$, $u_n\to U$ uniformly on $K$;
\item for every compact $K\subseteq\Omega$, $g_n\to G$ uniformly on $K$.
\end{enumerate}
Then $U$ is $p$-harmonic in $\Omega$ with weak gradient $G$.
\end{lemma}
\begin{proof}
We verify the six conditions in the definition of a $p$-harmonic function with weak gradient,
\noderef{IsPHarm}.  Throughout, "uniformly on $K$" has its usual meaning: for every
$\varepsilon>0$ there is $N$ such that for all $n\ge N$ and all $y\in K$ the relevant distance
is at most $\varepsilon$.

\medskip\noindent\emph{(1) $\Omega$ is open.}
Hypothesis (1) with $n=0$ gives that $u_0$ is $p$-harmonic in $\Omega$, and condition (1) of
\noderef{IsPHarm} for that datum says $\Omega$ is open.  In particular $\Omega$ is Lebesgue
measurable, and so is every compact subset of $\mathrm{Plane}$, by clause (2) of the background
proposition on integrability of continuous and compactly supported functions; this is used
repeatedly below, since continuity on a set yields almost everywhere strong measurability with
respect to the restricted measure only when the set is measurable.

\medskip\noindent\emph{(2) $U$ is continuous on $\Omega$.}
Let $x\in\Omega$.  Since $\Omega$ is open by step (1), there is $\delta>0$ with
$\overline{B_\delta(x)}\subseteq\Omega$.  Put $K:=\overline{B_\delta(x)}$, a compact subset of
$\Omega$.  By condition (2) of \noderef{IsPHarm}, each $u_n$ is continuous on $\Omega$, hence
on $K$.  By hypothesis (2) applied to $K$, $u_n\to U$ uniformly on $K$.  By the background
proposition on uniform limits of continuous functions, $U$ is continuous on $K$.  Since
$B_\delta(x)\subseteq K$ is a neighbourhood of $x$, the same proposition gives that $U$ is
continuous at $x$.  As $x\in\Omega$ was arbitrary, $U$ is continuous at every point of
$\Omega$, hence continuous on $\Omega$.

\medskip\noindent\emph{(3) $G$ is almost everywhere strongly measurable on $\Omega$.}
Let $y\in\Omega$.  The singleton $\{y\}$ is a compact subset of $\Omega$, so hypothesis (3)
applied to $K=\{y\}$ gives $g_n(y)\to G(y)$ (uniform convergence on a set implies pointwise
convergence at each of its points).  Thus $g_n\to G$ pointwise at every point of $\Omega$, and
therefore almost everywhere with respect to Lebesgue measure restricted to $\Omega$, which is
carried by $\Omega$.  By condition (3) of \noderef{IsPHarm} every $g_n$ is almost everywhere
strongly measurable with respect to that measure.  The background proposition on pointwise
limits and almost everywhere strong measurability now gives the same for $G$.

\medskip\noindent\emph{(4) For every compact $K\subseteq\Omega$, $\lvert G\rvert^{p}$ is
integrable on $K$.}
Fix such a $K$.  By hypothesis (3) applied to $K$ with $\varepsilon=1$ there is $N$ with
$\lvert g_n(y)-G(y)\rvert\le1$ for all $n\ge N$ and all $y\in K$; take $n=N$, so that by the
triangle inequality
\[
  \lvert G(y)\rvert\le\lvert g_N(y)\rvert+1\qquad\text{for every }y\in K .
\]
Since $p\ge1>0$, clause (3) of the background proposition on elementary properties of real
powers (monotonicity of $t\mapsto t^{p}$ on $[0,\infty)$) and then its clause (5) give
\[
  \lvert G(y)\rvert^{p}\le\bigl(\lvert g_N(y)\rvert+1\bigr)^{p}
  \le2^{p}\bigl(\lvert g_N(y)\rvert^{p}+1\bigr)\qquad(y\in K).
\]
By condition (4) of \noderef{IsPHarm} for the datum $n=N$, the function
$y\mapsto\lvert g_N(y)\rvert^{p}$ is integrable on $K$; the constant $1$ is integrable on $K$
because $K$ is compact and hence of finite measure.  So the right-hand side is integrable on
$K$.  The function $y\mapsto\lvert G(y)\rvert^{p}$ is measurable on $K$: $G$ is almost
everywhere strongly measurable by step (3), and composing with the continuous maps
$\eta\mapsto\lvert\eta\rvert$ and $t\mapsto t^{p}$ preserves that, by the composition clause of
the background proposition on pointwise limits.  A measurable function dominated by an
integrable function is integrable, by clause (3) of the integrability background proposition.
Hence $\lvert G\rvert^{p}$ is integrable on $K$.

\medskip\noindent\emph{Preparations for (5) and (6).}
Let $\varphi\colon\mathrm{Plane}\to\mathbb{R}$ be of class $C^\infty$ with compact support and
$\operatorname{supp}\varphi\subseteq\Omega$.  Put $K_0:=\operatorname{supp}\varphi$, a compact
subset of $\Omega$, and
\[
  B:=\sup_{y\in\mathrm{Plane}}\lvert D\varphi(y)\rvert<\infty,\qquad
  B':=\sup_{y\in\mathrm{Plane}}\lvert\varphi(y)\rvert<\infty,
\]
both finite because $\varphi$ and $D\varphi$ are continuous and vanish off the compact set
$K_0$ (on the open set $\mathrm{Plane}\setminus K_0$ the function $\varphi$ is identically $0$,
so all its derivatives vanish there).  Also $\lvert K_0\rvert<\infty$.

We record that the following six functions are integrable on $\Omega$ and vanish off $K_0$, so
that each of their integrals over $\Omega$ equals the corresponding integral over $K_0$ by
clause (4) of the integrability background proposition:
\[
  u_n\,D\varphi,\quad U\,D\varphi,\quad \varphi\,g_n,\quad \varphi\,G,\quad
  \bigl\langle A_p(g_n),D\varphi\bigr\rangle,\quad \bigl\langle A_p(G),D\varphi\bigr\rangle .
\]
Indeed:
\begin{itemize}
\item $u_n$ and $U$ are continuous on the measurable set $\Omega$ (condition (2) of
\noderef{IsPHarm} and step (2)) and $D\varphi$ is continuous on $\mathrm{Plane}$, so
$u_n\,D\varphi$ and $U\,D\varphi$ are continuous on $\Omega$, hence bounded and integrable on
the compact subset $K_0\subseteq\Omega$ by clause (3) of the integrability proposition; and they
vanish off $K_0$.
\item $g_n$ and $G$ are almost everywhere strongly measurable on $\Omega$ (condition (3) of
\noderef{IsPHarm} and step (3)), so $\varphi\,g_n$ and $\varphi\,G$ are too.  Taking $a=1$ in
clause (4) of the elementary-powers proposition, $\lvert\eta\rvert\le1+\lvert\eta\rvert^{p}$
for every vector $\eta$ (legitimate since $1\le p$), hence
\[
  \lvert\varphi(y)\,g_n(y)\rvert\le B'\bigl(1+\lvert g_n(y)\rvert^{p}\bigr),\qquad
  \lvert\varphi(y)\,G(y)\rvert\le B'\bigl(1+\lvert G(y)\rvert^{p}\bigr),
\]
and the right-hand sides are integrable on $K_0$ by condition (4) of \noderef{IsPHarm},
respectively step (4), together with $\lvert K_0\rvert<\infty$.  Clause (3) of the integrability
proposition then gives integrability on $K_0$, and both functions vanish off $K_0$.
\item By \noderef{PFluxContinuity} the flux $A_p$ is continuous on $\mathrm{Plane}$, so
$A_p\circ g_n$ and $A_p\circ G$ are almost everywhere strongly measurable on $\Omega$ by the
composition clause of the pointwise-limits proposition, and so are their inner products with the
continuous $D\varphi$.  By the norm identity of \noderef{PFluxContinuity} and Cauchy--Schwarz,
\[
  \bigl\lvert\bigl\langle A_p(g_n(y)),D\varphi(y)\bigr\rangle\bigr\rvert
  \le B\,\lvert A_p(g_n(y))\rvert=B\,\lvert g_n(y)\rvert^{\,p-1}
  \le B\bigl(1+\lvert g_n(y)\rvert^{p}\bigr),
\]
using clause (4) of the elementary-powers proposition with $a=p-1\in[0,p]$, and likewise with
$G$ in place of $g_n$.  Integrability on $K_0$ and vanishing off $K_0$ follow as before.
\end{itemize}

\medskip\noindent\emph{(5) $G$ is the distributional gradient of $U$ on $\Omega$.}
Let $\varphi$ be as above.  For every $n$, condition (5) of \noderef{IsPHarm} gives
\begin{equation}\label{eq:phl-n}
  \int_\Omega u_n\,D\varphi\,dy=-\int_\Omega \varphi\,g_n\,dy .
\end{equation}
Let $\varepsilon>0$.  By hypothesis (2) applied to $K_0$ there is $N_1$ with
$\lvert u_n(y)-U(y)\rvert\le\varepsilon$ for all $n\ge N_1$, $y\in K_0$; hence, for $n\ge N_1$,
\[
  \left\lvert\int_\Omega u_n\,D\varphi\,dy-\int_\Omega U\,D\varphi\,dy\right\rvert
  =\left\lvert\int_{K_0}(u_n-U)\,D\varphi\,dy\right\rvert
  \le\varepsilon\,B\,\lvert K_0\rvert .
\]
By hypothesis (3) applied to $K_0$ there is $N_2$ with
$\lvert g_n(y)-G(y)\rvert\le\varepsilon$ for all $n\ge N_2$, $y\in K_0$; hence, for $n\ge N_2$,
\[
  \left\lvert\int_\Omega \varphi\,g_n\,dy-\int_\Omega \varphi\,G\,dy\right\rvert
  =\left\lvert\int_{K_0}\varphi\,(g_n-G)\,dy\right\rvert
  \le\varepsilon\,B'\,\lvert K_0\rvert .
\]
As $\varepsilon>0$ was arbitrary and $B,B',\lvert K_0\rvert$ are finite, this says
\[
  \int_\Omega u_n\,D\varphi\,dy\longrightarrow\int_\Omega U\,D\varphi\,dy,
  \qquad
  \int_\Omega \varphi\,g_n\,dy\longrightarrow\int_\Omega \varphi\,G\,dy .
\]
Passing to the limit in \eqref{eq:phl-n}, and using that limits of convergent sequences are
unique and that negation is continuous,
\[
  \int_\Omega U\,D\varphi\,dy=-\int_\Omega \varphi\,G\,dy ,
\]
which is condition (5) for $U$ and $G$ on $\Omega$.

\medskip\noindent\emph{(6) $U$ solves the weak $p$-Laplace equation on $\Omega$.}
Let $\varphi$ be as above, and set
\[
  F_n(y):=\bigl\langle A_p(g_n(y)),D\varphi(y)\bigr\rangle,\qquad
  F(y):=\bigl\langle A_p(G(y)),D\varphi(y)\bigr\rangle .
\]
For every $n$, condition (6) of \noderef{IsPHarm} gives $\int_\Omega F_n\,dy=0$.

\emph{Pointwise convergence.}  Let $y\in\Omega$.  By step (3), $g_n(y)\to G(y)$.  By
\noderef{PFluxContinuity} the map $A_p$ is continuous on $\mathrm{Plane}$ --- and this is the
point at which the argument would fail if continuity at the origin were unavailable, since the
convergence may well take place at a critical point, where $G(y)=0$.  Hence
$A_p(g_n(y))\to A_p(G(y))$, and since the inner product is continuous in its first argument,
$F_n(y)\to F(y)$.  So $F_n\to F$ pointwise on $\Omega$, in particular almost everywhere with
respect to Lebesgue measure restricted to $\Omega$.

\emph{Domination.}  By hypothesis (3) applied to $K_0$ with $\varepsilon=1$ there is $N$ with
$\lvert g_n(y)-G(y)\rvert\le1$, hence $\lvert g_n(y)\rvert\le\lvert G(y)\rvert+1$, for all
$n\ge N$ and $y\in K_0$.  For such $n$ and $y$, the norm identity of \noderef{PFluxContinuity},
Cauchy--Schwarz, and clauses (3), (4) and (5) of the elementary-powers proposition give
\[
  \lvert F_n(y)\rvert
  \le B\,\lvert g_n(y)\rvert^{\,p-1}
  \le B\bigl(\lvert G(y)\rvert+1\bigr)^{\,p-1}
  \le B\Bigl(1+\bigl(\lvert G(y)\rvert+1\bigr)^{p}\Bigr)
  \le B\Bigl(1+2^{p}\bigl(\lvert G(y)\rvert^{p}+1\bigr)\Bigr)=:H(y),
\]
where monotonicity of $t\mapsto t^{\,p-1}$ (clause (3), applicable because $p-1>0$) was used in
the second step, clause (4) with $a=p-1\in[0,p]$ in the third, and clause (5) with $q=p\ge1$ in
the fourth.  Moreover $F_n$ vanishes off $K_0$, and so does $H$ once we replace $H$ by its
product with the indicator of $K_0$; with that convention $\lvert F_n\rvert\le H$ on all of
$\Omega$ for $n\ge N$.  The function $H$ is integrable on $\Omega$: it vanishes off $K_0$, and
on $K_0$ it is a finite linear combination of the constant $1$ and of
$\lvert G\rvert^{p}$, both integrable on $K_0$ by $\lvert K_0\rvert<\infty$ and step (4).

\emph{Conclusion.}  Each $F_n$ is almost everywhere strongly measurable with respect to
Lebesgue measure restricted to $\Omega$, as recorded in the preparations.  Discarding the
finitely many indices $n<N$ changes no limit, so the dominated convergence background
proposition applies to $(F_n)_{n\ge N}$ with dominating function $H$ and gives
\[
  \int_\Omega F_n\,dy\longrightarrow\int_\Omega F\,dy .
\]
Since $\int_\Omega F_n\,dy=0$ for every $n$, the limit is $0$, that is
\[
  \int_\Omega\bigl\langle A_p(G(y)),D\varphi(y)\bigr\rangle\,dy=0 ,
\]
which is condition (6) for $U$ and $G$ on $\Omega$.

\medskip\noindent All six conditions hold, so $U$ is $p$-harmonic in $\Omega$ with weak
gradient $G$.
\end{proof}
